\documentclass[tikz,border=4pt]{standalone}
\usepackage{ctex}
\usepackage{pgfplots}
\pgfplotsset{compat=1.18}
\usetikzlibrary{calc}

\begin{document}
% part11/04 方差大小对应数据分布散度示意图
% 左：小方差（集中），右：大方差（分散）
\begin{tikzpicture}[scale=1.0, thick]
  % ===== 左图：小方差，数据集中 =====
  \begin{scope}[xshift=0cm]
    \node[font=\small, above] at (5,4.0) {小方差（$S^2$ 小）};
    \node[font=\footnotesize, above] at (5,3.6) {数据集中在均值附近};
    \draw[->] (2,0) -- (8,0) node[right, font=\small] {数据值};
    \draw[->] (2,0) -- (2,3.5) node[above, font=\small] {频率};
    % 均值竖线
    \draw[red, dashed] (5,0) -- (5,3.5) node[above, red, font=\small] {$\bar{x}$};
    % 近似正态曲线（集中，窄高）
    \draw[blue, thick, domain=3:7, samples=60]
      plot (\x, {3.0*exp(-3*(\x-5)^2)});
  \end{scope}

  % ===== 右图：大方差，数据分散 =====
  \begin{scope}[xshift=8cm]
    \node[font=\small, above] at (5,4.0) {大方差（$S^2$ 大）};
    \node[font=\footnotesize, above] at (5,3.6) {数据分散在均值两侧};
    \draw[->] (2,0) -- (8,0) node[right, font=\small] {数据值};
    \draw[->] (2,0) -- (2,3.5) node[above, font=\small] {频率};
    % 均值竖线
    \draw[red, dashed] (5,0) -- (5,3.5) node[above, red, font=\small] {$\bar{x}$};
    % 近似正态曲线（分散，宽矮）
    \draw[blue, thick, domain=2.5:7.5, samples=60]
      plot (\x, {1.1*exp(-0.5*(\x-5)^2)});
  \end{scope}
\end{tikzpicture}
\end{document}
